% 復習2　問いの立て方（記入例）
\session{2}{1}{60分}{問いの立て方}{AIとブレスト}{}

\goals{
  \item テーマとリサーチ・クエスチョン（問い）の違いを説明できる
  \item 良い問いの条件（調べて答えられる／範囲が適切／答える意義がある）を挙げられる
  \item AIとのブレストで視点を広げつつ、最後は自分で選ぶ
}

\work{確認}{テーマと問いの違い}
\begin{wstab}{colspec={Q[l,wd=40mm,m] X[l,m]},row{1-Z}={ht=9mm}}
  \hc{例：テーマ} & 小売業のキャッシュレス決済 \\
  \hc{例：問い} & なぜ小規模な飲食店では、キャッシュレス決済の導入が遅れているのか \\
  \hc{自分のテーマ} & \af{中小企業と生成AI} \\
  \hc{問いにすると} & \af{愛知県の中小企業では、なぜ生成AIの業務利用が進まないのか} \\
  \hc{良い問いの3条件} & ①\,\af{調べて答えられる}\hspace{8mm} ②\,\af{範囲が広すぎず狭すぎない}\hspace{8mm} ③\,\af{答える意義がある} \\
\end{wstab}

\work[60mm]{ワーク1}{AIで視点を広げる}
\begin{promptbox}[視点を広げる]
私は大学生で、「（あなたのテーマ）」に関心があります。このテーマについて、経営戦略、組織、マーケティング、情報システム、経済、社会の6つの視点から、それぞれ2つずつ研究の問いを提案してください。各問いは「なぜ〜なのか」「〜は〜にどのように影響するのか」の形で、1行ずつ箇条書きにしてください。
\end{promptbox}
\instr{AIの候補のうち、気になったものを視点ごとに書き写す（言い換えてよい）。}
\begin{wstabh}{colspec={Q[l,wd=22mm,m] X[l,m]},row{2-Z}={ht=10mm}}
  視点 & 気になった問いの候補 \\
  経営戦略 & \af{生成AIの導入は、中小企業の競争力にどのように影響するのか} \\
  組織 & \af{なぜ中小企業では、生成AIを使いこなす人材が育ちにくいのか} \\
  マーケティング & \af{生成AIは、中小企業の販促や顧客対応をどのように変えるのか} \\
  情報システム & \af{既存の業務システムの状況は、生成AIの導入にどう影響するのか} \\
  経済 & \af{導入の費用と効果は、中小企業の導入判断にどう影響するのか} \\
  社会 & \af{情報漏えいへの不安は、中小企業の生成AI利用をどのように妨げているのか} \\
\end{wstabh}

\work[60mm]{ワーク2}{3条件で採点し、上位3つに絞る}
\instr{各条件を 3（十分）・2（やや）・1（不十分）で採点する。必要ならAIに評価と狭める案を求める。}
\begin{promptbox}[問いを絞る]
次の問いについて、「大学生が数か月で調べて答えられるか」「範囲は適切か」「答えることに意義があるか」の3点から評価し、必要なら狭める案を示してください。問い：（候補の問い）
\end{promptbox}
\begin{wstabh}{colspec={X[l,m] Q[c,wd=17mm,m] Q[c,wd=17mm,m] Q[c,wd=17mm,m] Q[c,wd=12mm,m] Q[c,wd=10mm,m]},row{2-Z}={ht=10mm}}
  問いの候補 & 調べて\newline 答えられる & 範囲が\newline 適切 & 答える\newline 意義 & 合計 & 順位 \\
  \af{なぜ中小企業では生成AIの業務利用が進まないのか} & \ans{3} & \ans{2} & \ans{3} & \ans{8} & \ans{1} \\
  \af{情報漏えいへの不安は、生成AI利用をどう妨げているか} & \ans{2} & \ans{3} & \ans{3} & \ans{8} & \ans{2} \\
  \af{生成AIを使いこなす人材が育ちにくいのはなぜか} & \ans{2} & \ans{2} & \ans{3} & \ans{7} & \ans{3} \\
  \af{生成AIの導入は中小企業の競争力にどう影響するか} & \ans{1} & \ans{1} & \ans{3} & \ans{5} & \\
  \af{導入の費用と効果は、導入判断にどう影響するか} & \ans{1} & \ans{2} & \ans{2} & \ans{5} & \\
\end{wstabh}

\work{セルフチェック}{「それはどうやって調べるの？」に答える}
\instr{上位の問いについて、AIに「この問いを調べる方法について、学生に鋭い質問を3つしてください」と頼み、答えてみる。友人と説明し合ってもよい。}
\atwoboxes{受けた質問}{自分の答え・気づいたこと}{30mm}
{「進まない」とは、何と比べて進んでいないのか？ 大企業とか、他の地域とか？\newline 「業務利用」に、社員が個人的に使う場合は含むのか？}
{白書の統計で、企業規模別の利用状況を比べる。「業務利用」は会社として業務に使っている場合に限る。比べる相手と言葉の意味を決めないと調べられないと気づいた。}

\work{決定}{仮のリサーチ・クエスチョン}
\abox{18mm}{愛知県の中小企業では、なぜ生成AIの業務利用が進まないのか。}

\begin{nexttime}
  \item 仮の問いを選んだ理由を200字程度でまとめる
  \item 使ったプロンプトと、AIの案のうち採用しなかったものとその理由を、下のAI活用記録に書く
\end{nexttime}
\abox[選んだ理由（200字程度）]{40mm}{アルバイト先の町工場の社長が、生成AIを使ってみたいが何に使えばよいかわからないと話していた。生成AIの話題は大企業の事例が多いが、地域の経済を支える中小企業での状況は身近なのに知られていないと感じた。公的な白書や統計で全体の傾向を確かめつつ、身近な企業に話を聞くことで、数か月で答えに近づける範囲だと考えた。導入に悩む企業にとっても、答える意義がある問いだと考えた。}
\aimemoA{問いの候補出し（6つの視点）と、候補の採点}
{「6つの視点から問いを2つずつ」「3条件で評価し、必要なら狭める案を示して」}
{12の候補から6つを書き出した。「競争力への影響」は、効果を測るのが難しく範囲も広いため不採用。採点はAIの評価を参考にしつつ、意義の点数は自分で付け直した。}
{「中小企業で生成AIの利用が少ない」という前提が、白書などで実際に示されているかを確かめる予定（復習3）}

\begin{point}
  \item テーマ（話題の範囲）と問い（答えを求める疑問文）は違う。「〜について」は問いではない。
  \item 良い問いの3条件は「調べて答えられる」「範囲が適切」「答える意義がある」の3つ（模範解答）。
  \item 採点の数値そのものより、点数の理由を説明できるかが大切。AIの評価をそのまま写さない。
  \item 「なぜ進まないのか」のような問いは、「何と比べて」「どの範囲で」を決めると調べられる形になる。セルフチェックの質問は、そのための手がかり。
  \item AI活用記録には「採用しなかった案と理由」を必ず書く。
\end{point}
